% 致谢
时间的车轮滚滚向前，推着我走过了七年求学时光。当名为学生生涯的这本书翻到终章，细数其中写满的青涩与汗水，那个懵懂的少年即将告别校园，踏入人生的新篇章，不禁感慨万千。

回首这段求学之路，我首先要感谢我的导师卢昱杰教授。卢老师是我本科时期的班主任，正是他对于学生职业发展与身心健康的真诚关怀以及丰富的学识与眼界，吸引我加入了课题组。在科研工作中，卢老师对学生的求学态度和工作质量严格要求、精益求精。得益于此，我在进组后的极短时间内，便在逻辑思维和问题解决能力上得到了全面提升。卢老师的高标准与严要求，是我在这个日新月异的时代里最宝贵的财富，时刻提醒我以严谨之心治学，以敬畏之心律己，不断拓宽认知的边界与视野的广度。

同窗之谊，弥足珍贵。感谢遇见课题组的小伙伴们，在这条披荆斩棘的路上，你们给我的帮助与体谅，远非三言两语能道尽。感谢王硕、杨川、雁杰、仲涛、王娜、瑞涵，在科研与生活的方方面面给予我无私的指导，你们永远是我心中最可靠的前辈。感谢孙一玺在我毕业论文数据采集工作上的鼎力相助，可以说，没有你就没有这篇毕业论文。感谢同期的利剑、修龙、缘诚，我们是并肩前行的队友，一路共同披荆斩棘、共同成长。曲终未必人散，愿我们友谊长存。

纸短情长，言浅恩深。感谢一直在我身后的家人和朋友，研究生阶段经历了许多意想不到的波折，没有你们的支持与陪伴，我无法度过那无数个迷惘的夜晚。正是这些经历让我逐渐明白，你们是我最宝贵的财富。我爱你们，愿你们身体健康、平安喜乐。

最后，也感谢一下一路坚持的自己吧。站在旅途的终点，会有一种奇妙的视角，每每回望来时路，都会惊叹于自己如何一步步成为今天的模样。那些曾经以为过不去的坎，都化作了人生路上的一面面旗帜，随着越走越远，显得越发渺小了。而那些没能做好的地方，如果重来一次，人生又会怎样？我不知道。我只知道，人生没有如果，只有勇往直前。那就往前走吧，虽然不知路在何方，虽然内心仍会迷惘，但只要不停下脚步，这条平凡之路，便不会平凡。

心中纵有千言万语，行文至此终要作别。自此奔赴人生旷野，愿我们都得偿所愿，成为自己想成为的人。
